
\begin{proposition}[Proved actual occurrence of both Hopf types]\label{cdp:st:C.3.6}

There are compact complex threefolds with $b_1=b_2=0$ and $a=1$
whose holomorphic algebraic reduction is a smooth bundle over
$\PP^1$ with primary Hopf fibre. There are also such threefolds
with secondary Hopf fibre. In these examples every fibre has Euler
number zero and $c_3=0$.
\end{proposition}
\begin{proof}
Put $a_1=1$, $a_2=\sqrt2$ and
$W=(\C^2\setminus\{0\})_z\times(\C^2\setminus\{0\})_w$.
The additive complex Lie group $\C$ acts holomorphically by
\[
t\cdot(z,w_1,w_2)
 =\bigl(e^t z,e^{ia_1t}w_1,e^{ia_2t}w_2\bigr).
\tag{C.H.1}\label{cdp:eq:C.H.1}
\]
This action is free. Indeed $e^t z=z$ implies $t=2\pi in$ with
$n\in\Z$; a nonzero $w_j$ then gives $e^{-2\pi a_jn}=1$ and
therefore $n=0$. It is proper as well. If a point and its translate
range in fixed compact subsets of $W$, the $z$-norm bounds
$\operatorname{Re}t$. Writing $s=\operatorname{Im}t$, the identity
\[
\|t\cdot w\|^2
 =e^{-2a_1s}|w_1|^2+e^{-2a_2s}|w_2|^2
\]
and the positive constants $a_1,a_2$ bound $s$ both above and
below, since both $w$-norms are bounded away from zero. The inverse
image of a compact subset under the action graph is consequently
closed and bounded in the group variable and compact in the other
variable. This proves properness.

Every orbit meets the real submanifold
$\{\|z\|=\|w\|=1\}=S^3\times S^3$ exactly once.
The first equation fixes $\operatorname{Re}t=-\log\|z\|$.
The displayed squared $w$-norm is strictly decreasing in $s$,
with limits $+\infty$ and $0$, so the second equation has one
solution. Its derivative is strictly negative, and the implicit
function theorem makes the resulting cross-section smooth.
Thus the complex quotient $Y=W/\C$ is a smooth compact complex
threefold diffeomorphic to $S^3\times S^3$. In particular
$b_1(Y)=b_2(Y)=0$ and $e(Y)=0$.

The map
\[
f:Y\longrightarrow\PP^1,\qquad [z,w]\longmapsto[z_1:z_2]
\]
is well-defined and holomorphic. Choose a holomorphic nonzero
frame of the tautological line on a small base disc. Scaling the
$z$-coordinate to this frame leaves exactly the transformations
$t=2\pi in$, which act on $w$ by powers of
\[
Q=\operatorname{diag}(q_1,q_2),\qquad
q_j=e^{-2\pi a_j}.
\]
This identifies $f$ over that disc with the product having fibre
$H=(\C^2\setminus\{0\})/\langle Q\rangle$.
The quotient is a primary Hopf surface, and these local products
prove that $f$ is a proper holomorphic submersion with connected
fibres.

We verify $a(H)=0$ directly. A meromorphic function on $H$ lifts
to a $Q$-invariant meromorphic function on $\C^2\setminus\{0\}$.
It extends meromorphically across the origin: the closure of its
polar divisor is analytic by the codimension-two extension theorem
for divisors, is locally principal on the smooth ball, and
multiplying by a local defining function reduces extension to
Hartogs' theorem for a holomorphic function. Write the germ as
$g/h$ with relatively prime convergent power series. The identity
$(g\circ Q)h=g(h\circ Q)$ gives
\[
g\circ Q=ug,\qquad h\circ Q=uh
\]
for one holomorphic unit $u$: divisibility follows from relative
primality, and preservation of the vanishing order makes the
quotient a unit. Write $u/u(0)=e^\ell$ with $\ell(0)=0$ and set
\[
v(w)=\exp\left(\sum_{n\ge0}\ell(Q^nw)\right).
\]
The sum converges normally on a sufficiently small ball, because
$Q$ is a strict contraction. It gives
$v(Qw)/v(w)=u(0)/u(w)$. Consequently both $vg$ and $vh$ are
eigenfunctions of pullback by $Q$ with the same eigenvalue $u(0)$.
The monomial $w_1^rw_2^s$ has eigenvalue
$e^{-2\pi(r+s\sqrt2)}$. These numbers are distinct for distinct
$(r,s)\in\Z_{\ge0}^2$. Comparing power-series coefficients shows
that each nonzero eigenfunction is a single monomial and that the
two eigenfunctions are proportional. The meromorphic function is
therefore constant, first near the origin and then everywhere by
analytic continuation.

Every meromorphic function on $Y$ is constant on each general
fibre of $f$. On a product chart, restricting to a local section
gives the descended meromorphic function on the base; equality on
the fibres and the identity theorem identify its pullback with
the original function. These local functions agree on overlaps
and extend at the finitely many possible poles on the base.
Thus $\mathcal M(Y)=f^*\mathcal M(\PP^1)$, proving $a(Y)=1$
and proving that $f$ is the holomorphic algebraic reduction.

For any integer $m>1$, let $\mu_m$ act on $W$ by
$(z,w)\mapsto(z,\zeta w)$. It commutes with \eqref{cdp:eq:C.H.1} and induces
a free action on $Y$. A fixed orbit would give $t=2\pi in$ from
the $z$-coordinate; equality of norms in a nonzero $w_j$ then
gives $n=0$ and $\zeta=1$. The quotient
$Y_m=Y/\mu_m$ is a smooth compact complex threefold. Transfer
injects $H^j(Y_m;\Q)$ into $H^j(Y;\Q)$, so $b_1(Y_m)=b_2(Y_m)=0$.
The simply connected cover $Y\simeq S^3\times S^3$ gives
$\pi_1(Y_m)=\mu_m$, and the universal coefficient theorem gives
$H^2(Y_m;\Z)\simeq\Z/m\Z$: its free part is zero by the
rational calculation, and its torsion is
$\operatorname{Ext}(H_1(Y_m;\Z),\Z)$. Thus these secondary
examples meet exactly the rational hypotheses of C.3; the primary
example also has $H^1(Y;\Z)=H^2(Y;\Z)=0$.
The induced map to $\PP^1$ has fibre
\[
H_m=(\C^2\setminus\{0\})/
       \bigl(\langle Q\rangle\times\mu_m\bigr).
\]
Its universal cover is $\C^2\setminus\{0\}$ and its fundamental
group is $\Z\times\mu_m$, which is noncyclic. It is therefore a
secondary Hopf surface, with the displayed primary cover of degree
$m$. Any meromorphic function on $H_m$ pulls back to one on $H$,
so $a(H_m)=0$. The same local descent proves
$\mathcal M(Y_m)=f_m^*\mathcal M(\PP^1)$ and $a(Y_m)=1$.
Finally transfer of Euler characteristic under the finite covering,
or the smooth-fibration Euler formula, gives $e(Y_m)=0$.
\end{proof}


